\subsubsection{DS-forskning (Hevner)}
Valet av metod grundade sig i att DS-forskning vilar på en innovativ grund, och paradigmen med järnvägstrafik med C-DAS-stöd är ung och på många ställen ej ännu införd.
Innovationsfokuset skulle kunna lämpa sig väl för kontexten C-DAS då det inte är säkert att de mest optimala lösningar ännu framtagits eller täcks av befintlig teknik.
Design-science research drivs av relevans och trovärdighet vilket passar problem- och frågeställningen bra.
Relevansen fastställs i att verksamhetsproblemet är tydligt utstakat och väl förankrat i verkligheten.
Trovärdigheten täcks av att datautbytesprotokollet utvärderas, valideras och förfinas.
\subsubsection{Fallstudie (Oates)}
En fallstudie såsom definierat enligt Oates lämpar sig väl för djuplodande analys av särskilt specifika problem~\cite{Oates2005}.
Kvalitativ datainsamling i form av semi-strukturella intervjuer samt dokumentanalys såsom de tillämpats inom ramarna för denna studie skulle kunna lämpa sig väl även i en fallstudie.
Problemet är att en fallstudie i grunden tar fasta på redan existerande system eller fenomen.
C-DAS är inte fullt utrullat i många länder idag vilket skulle kunna ses som mindre gynnsamma förutsättningar för att genomföra en fallstudie.
En fallstudie skulle inte nödvändigtvis mynna ut i innovativa förslag på framtagande eller, i fallet för den här studien, förbättringspotential av det innevarande protokoll från lokförares perspektiv.
%För fallet med C-DAS är systemet inte vida etablerat varför det skulle kunna argumenteras för att fallstudie inte vore optimal forskningsmetod.